\section{Government Agent：从 Task Demo 到 Workflow System}

全讲最后一个问题最接近今天的 agent engineering。Alswaha 区分 legacy B2B SaaS 的 AI enhancement 与 AI-native player，并认为难点不只在模型，而在 data、business model 和 human-in-the-loop。更精确地说，task-level success 只有跨过 workflow integration、multi-agent coordination 和 public accountability 三道门，才能成为公共服务系统。

\subsection{Task、workflow 与 multi-agent 的三次状态扩张}

\term{task-level agent} 在局部输入上完成 drafting、classification 或 lookup；\term{workflow-level agent} 跨多个系统、权限和长状态推进 case；\term{multi-agent system} 让多个 agent 分工、共享 memory、协调工具和处理冲突。每次扩张都会增加 state space、failure interaction 与 audit complexity。

\lecturefigure{16-government-agent-maturity.png}{政府 Agent 要依次跨越 task、workflow、multi-agent 和 public accountability。}{本地字幕 49:50--52:17；概念重绘。}

\begin{knowledgebox}{读图：为什么 task demo 最容易成功}
Task demo 的输入可人工筛选、状态短、权限少、失败可忽略；真实 workflow 会遇到缺失字段、重复 case、跨部门 ownership、deadline、appeal 和 policy exception。Multi-agent 又引入 shared memory、conflicting plans、deadlock 与 cascading tool error，因此不能用单轮 benchmark 外推生产成功率。
\end{knowledgebox}

\subsection{Data layer、human-in-the-loop 与 business model}

\term{human-in-the-loop}（人在回路）不是“任何结果都让人点确认”，而是把高风险、低置信、政策例外和申诉路径交给有权限的人，并收集其理由作为改进证据。Data layer 需要 canonical entity、lineage、quality rule、permission 与 retention；business model 则回答谁为部署、错误、人工复核和持续更新付费。

\begin{importantbox}{公共服务 Agent 的验收矩阵}
\begin{tabularx}{\textwidth}{L{0.18\textwidth}L{0.29\textwidth}X}
\toprule
维度 & 关键问题 & 可测证据 \\
\midrule
任务质量 & 输出是否正确、完整、可解释？ & expert review、error taxonomy、calibration \\
流程完整性 & 是否正确推进 state 和 exception？ & end-to-end case success、rework、deadline miss \\
公平与权利 & 是否可申诉、是否对群体产生差异影响？ & subgroup audit、appeal outcome、override reason \\
安全与隐私 & 权限、数据和 tool 是否最小化？ & access log、red-team、leakage/abuse test \\
经济性 & 自动化是否真的减少 total cost / cycle time？ & compute + integration + human review + incident cost \\
\bottomrule
\end{tabularx}
\end{importantbox}

\begin{warningbox}{“模型不是难点”也不能绝对化}
Data 和 workflow 常是主要瓶颈，但 base model 的 language coverage、reasoning、tool reliability、calibration 与 robustness 仍会限制系统。更准确的结论是：模型能力是必要条件，生产成功由 model、data、process、people、policy 和 economics 的乘积决定。
\end{warningbox}

\teachervoice{全讲最值得迁移的 teacher voice 出现在结尾：一些团队把 vertical data 直接接到模型，只得到 task augmentation；真正的 AI-native 系统会重新设计 data creation、human feedback 和 workflow ownership。}

\subsection{本章小结}

Government Agent 的边界不在聊天界面，而在 case state、权限、数据、例外、申诉和责任。Task demo 可以由模型驱动；workflow system 必须由完整的社会技术架构驱动。

